\chapter{Aggregation Pipeline Mastery II: Window Functions and Materialized Views}
\index{\$setWindowFields}\index{window functions}\index{moving average}\index{\$denseRank}\index{\$derivative}\index{materialized view}\index{\$merge}\index{on-demand materialized view}%

\begin{objectives}
\begin{itemize}
    \item Compute running totals, rankings, derivatives, and moving averages with \texttt{\$setWindowFields} partitions and window bounds.
    \item Distinguish accumulator behavior across \texttt{\$group}, \texttt{\$project}, and \texttt{\$setWindowFields} contexts.
    \item Respect the 100 MB per-stage memory boundary and the sort requirements of windowed computation.
    \item Build incrementally refreshed materialized views with \texttt{\$merge} and deterministic keys.
    \item Deliver a 30-day moving-average financial analytics pipeline in the Thursday lab.
    \item Architect a real-time leaderboard over fifty million gaming events.
\end{itemize}
\end{objectives}

\begin{crosscloud}
Window functions are the most portable analytical concept in this book: SQL \texttt{OVER (PARTITION BY ... ORDER BY ...)} exists verbatim in Redshift, Synapse, BigQuery, and Snowflake. MongoDB's \texttt{\$setWindowFields} is the same mathematics with document semantics.

\vspace{2mm}
{\footnotesize
\begin{tabular}{@{}p{2.9cm}p{3.0cm}p{3.0cm}p{3.0cm}@{}} 
\toprule
\textbf{Concept} & \textbf{AWS (Redshift)} & \textbf{Azure (Synapse)} & \textbf{GCP (BigQuery)} \\
\midrule
Window stage & \texttt{OVER()} clause & \texttt{OVER()} clause & \texttt{OVER()} clause \\
Rank & \texttt{RANK}/\texttt{DENSE\_RANK} & Same & Same \\
Moving average & \texttt{AVG() OVER (ROWS)} & Same & Same \\
Derivative & \texttt{LAG()} arithmetic & Same & Same \\
Materialized view & Redshift MVs (auto-refresh) & Indexed views / Spark SQL & BigQuery MVs (auto-refresh) \\
Incremental refresh & Automatic & Manual / pipeline & Automatic, partition-aware \\
\addlinespace
Snowflake & \multicolumn{3}{l}{Window functions identical; materialized views auto-maintained; streams + tasks for incremental pipelines} \\
\bottomrule
\end{tabular}}

\vspace{1mm}
MongoDB's distinction is the manual refresh path: \texttt{\$merge} into a results collection is a materialized view you maintain yourself, which means you also own its freshness, idempotency, and backfill story \cite{mongodbSetWindowFieldsDocs,kleppmann2017ddia}.
\end{crosscloud}

\section{Week 11 Roadmap: From Aggregates to Windows}
Chapter 10's \texttt{\$group} collapses many documents into one summary row. Windows answer the complementary question: compute an aggregate \emph{for each} document from its neighbors---the running balance at every ledger entry, each player's rank within their region, the 30-day moving average beside every transaction. This week also closes the analytics loop: windowed results that dashboards recompute per request should become materialized views, refreshed incrementally and provably \cite{mongodbSetWindowFieldsDocs}.

\begin{definitionbox}
A \textbf{window function} computes a value for each document from a \emph{frame} of related documents defined by a partition key, a sort order, and window bounds, without collapsing the stream the way \texttt{\$group} does.
\end{definitionbox}

\begin{keyterms}
\textbf{Partition}: the group of documents over which one window computes (e.g. one account). \textbf{Frame}: the subset of the partition feeding one document's computation (\texttt{range: [-30, 0], unit: "day"}). \textbf{\$denseRank}: rank without gaps. \textbf{\$derivative}: rate of change between frame endpoints. \textbf{Materialized view}: a collection holding precomputed pipeline output, refreshed by \texttt{\$merge}.
\end{keyterms}

\section{Monday: \texttt{\$setWindowFields} in Depth}
\subsection{Partitions, Sort, and Frames}
The stage takes three structural decisions. \texttt{partitionBy} groups the stream (one account, one player, one device). \texttt{sortBy} orders within the partition---mandatory for any frame-based operator, because a window without order is nondeterministic. The \texttt{window} bounds pick the frame: \texttt{documents: [-10, 0]} (last ten rows), \texttt{range: [-30, 0], unit: "day"} (last thirty days), or \texttt{documents: ["unbounded", "current"]} (running total).

\begin{lstlisting}[language=JavaScript, caption={Running balance and 30-day moving average per account.}]
db.ledger.aggregate([
  { $match: { postedAt: { $gte: start } } },
  { $setWindowFields: {
      partitionBy: "$accountId",
      sortBy: { postedAt: 1 },
      output: {
        runningBalance: { $sum: "$amount",
          window: { documents: ["unbounded", "current"] } },
        avg30d: { $avg: "$amount",
          window: { range: [-30, 0], unit: "day" } }
      } } }
]);
\end{lstlisting}

\subsection{Ranking and Change Operators}
\texttt{\$rank}, \texttt{\$denseRank}, and \texttt{\$documentNumber} order documents within partitions; \texttt{\$denseRank} leaves no gaps after ties, which is what leaderboards usually want. \texttt{\$derivative} computes the rate of change between the frame's endpoints---queue depth per minute, temperature drift per hour---and \texttt{\$expMovingAvg} variants weight recent values more heavily for smoothing noisy telemetry.

\begin{pitfall}
\textbf{Windows without deterministic sort.} If \texttt{sortBy} has ties and the frame is document-bounded, ties can fall on either side of a frame boundary between runs. Add a tiebreaker (\texttt{sortBy: \{ ts: 1, \_id: 1 \}}) whenever frame membership affects results.
\end{pitfall}

\subsection{Memory and the Sort Requirement}
\texttt{\$setWindowFields} must materialize each partition in sorted order; it is a blocking stage subject to the 100 MB limit and \texttt{allowDiskUse}. Two design consequences follow. First, filter before the window: partitioning a hundred-million-document stream you will later discard is pure waste. Second, the stage can exploit an index whose order matches \texttt{partitionBy} + \texttt{sortBy}, avoiding the explicit sort---the window analogue of Chapter 9's sort elimination.

\section{Tuesday: Accumulators Across Contexts}
The same operator name behaves differently by context, and confusing them is the classic interview trap. In \texttt{\$group}, \texttt{\$sum: "\$amount"} accumulates across the group and outputs one value. In \texttt{\$project}, \texttt{\$sum} with an array argument reduces that array within one document. In \texttt{\$setWindowFields}, the accumulator computes over the frame and outputs one value \emph{per document}. A pipeline review should always ask: which context is this accumulator in, and what is its input stream?

\begin{tipbox}
When a pipeline needs both per-group totals and per-document windows, compute the window first (it preserves every document), then \texttt{\$group}. The reverse order loses the documents the window needed.
\end{tipbox}

\section{Wednesday: Materialized Views with \texttt{\$merge}}
\subsection{The Pattern}
MongoDB has no automatically maintained materialized view; the idiom is an aggregation pipeline ending in \texttt{\$merge}, which upserts pipeline output into a target collection \cite{mongodbMergeDocs}. Run on a schedule---or triggered by change streams (Chapter 16)---it produces a read-optimized results collection that dashboards query as plain documents. Because you own the refresh, the pattern's three non-negotiables are deterministic keys (so re-runs upsert rather than duplicate), a refresh watermark (so runs are incremental), and idempotency (so replays are safe).

\begin{lstlisting}[language=JavaScript, caption={Incremental materialized view refresh with a watermark.}]
const wm = db.mv_state.findOne({ _id: "acct_metrics" })?.watermark
         ?? ISODate("2000-01-01");
const now = new Date();

db.ledger.aggregate([
  { $match: { postedAt: { $gte: wm, $lt: now } } },       // delta window
  { $setWindowFields: {
      partitionBy: "$accountId", sortBy: { postedAt: 1, _id: 1 },
      output: { avg30d: { $avg: "$amount",
        window: { range: [-30, 0], unit: "day" } } } } },
  { $merge: { into: "acct_metrics", on: "_id",
              whenMatched: "replace", whenNotMatched: "insert" } }
]);

db.mv_state.updateOne({ _id: "acct_metrics" },
  { $set: { watermark: now } }, { upsert: true });
\end{lstlisting}

\begin{notebox}
Windows look backwards, so a naive watermark misses late-arriving rows that change old frames. Production refreshes re-compute a lookback window (e.g. watermark minus 31 days for a 30-day window) so stragglers are absorbed by the next run.
\end{notebox}

\subsection{Correctness Contract}
The contract mirrors Chapter 10's: a reference implementation (pandas or SQL) validates a sampled window, the merge key is deterministic, and a reconcile job periodically recomputes a full partition and diffs it against the materialized collection. A materialized view without a reconcile job is a rumor with a schema.

\section{Thursday Hands-On Project: Financial Moving-Average Pipeline}
Build the transaction-log analytics lab: two million ledger entries, a 30-day moving average and exponential smoothing per account, and an incremental \texttt{\$merge} refresh.

\subsection{Data and Window Pipeline}
\begin{lstlisting}[language=JavaScript, caption={Ledger generator and windowed analytics.}]
use windowLab;
db.ledger.drop();
const batch = [];
for (let i = 0; i < 2000000; i++) {
  batch.push({ _id: i, accountId: "acct-" + (i % 50000),
    postedAt: new Date(Date.now() - (i % 90) * 86400000 - (i % 86400) * 1000),
    amount: ((i % 20000) / 100 - 100) });
  if (batch.length === 10000) {
    db.ledger.insertMany(batch, { ordered: false }); batch.length = 0;
  }
}
db.ledger.createIndex({ accountId: 1, postedAt: 1 });

db.ledger.aggregate([
  { $match: { accountId: { $in: ["acct-1", "acct-2"] } } },
  { $setWindowFields: {
      partitionBy: "$accountId", sortBy: { postedAt: 1, _id: 1 },
      output: {
        avg30d: { $avg: "$amount",
          window: { range: [-30, 0], unit: "day" } },
        ema: { $expMovingAvg: { input: "$amount", alpha: 0.1 } }
      } } }
]);
\end{lstlisting}

\subsection{Exercises Within the Lab}
Verify determinism by running the pipeline twice and diffing outputs; then remove the \texttt{\_id} tiebreaker, re-run, and observe which frames changed. Finish by wiring the incremental \texttt{\$merge} refresh from Wednesday and re-running it after inserting a late-arriving transaction dated 10 days back---the lookback window must absorb it.

\subsection*{Deliverables}
\begin{itemize}
    \item The generator, window pipeline, and incremental refresh script.
    \item The determinism experiment (with and without the sort tiebreaker).
    \item The late-arrival test showing the lookback window correcting the materialized view.
\end{itemize}

\subsection*{Acceptance Criteria}
\begin{itemize}
    \item Moving-average values match a pandas reference within floating tolerance on a sampled account.
    \item The refresh is idempotent: three consecutive runs produce identical \texttt{acct\_metrics} content.
    \item No stage spills with \texttt{allowDiskUse} disabled on the filtered stream.
\end{itemize}

\subsection*{Stretch Goals}
\begin{itemize}
    \item Re-implement the same metrics with Spark or DuckDB and compare runtime, memory, and operational complexity.
    \item Trigger the refresh from a change stream on \texttt{ledger} instead of a cron schedule.
\end{itemize}

\section{Friday Use Case and Architecture: Real-Time Gaming Leaderboard}
\subsection{Scenario and Requirements}
A gaming platform streams fifty million match events daily and must serve per-player global and regional leaderboards with sub-second reads, plus per-player rolling statistics (win rate over the last 20 matches, score trend). Event volume rules out computing windows at read time; the leaderboard must be materialized and incrementally maintained.

\subsection{Architecture}
The design separates the hot write path from the analytical refresh. Match events land in a sharded \texttt{events} collection (Chapter 10). A scheduled micro-batch pipeline---every minute---computes \texttt{\$setWindowFields} over new events partitioned by player, merges rank and rolling-stat documents into \texttt{leaderboard} and \texttt{player\_stats} collections keyed by player and window, and serves reads from those precomputed collections with plain indexed lookups. Regional boards are one \texttt{partitionBy} change, not a second system. Chapter 15 replaces the cron with a change-stream trigger once CDC is in the toolkit.

\begin{pitfall}
\textbf{Ranking fifty million players in one window.} Global rank is a single partition---one giant sort. Compute ranks per region (bounded partitions) and derive the global board by merging ranked regional streams, or accept an approximate global board updated less frequently.
\end{pitfall}

\section{Interview Q\&A}
\subsection{Windows and Views}
\begin{enumerate}
    \item \textbf{Q: \texttt{\$rank} versus \texttt{\$denseRank}?}\\
    A: After a tie, \texttt{\$rank} skips positions (1, 2, 2, 4); \texttt{\$denseRank} does not (1, 2, 2, 3). Leaderboards usually want dense ranks.

    \item \textbf{Q: Why does \texttt{\$setWindowFields} require \texttt{sortBy}?}\\
    A: Frame membership is defined by order; without a sort, ``the last 30 days per row'' is meaningless and results are nondeterministic.

    \item \textbf{Q: How does a MongoDB materialized view differ from Snowflake's?}\\
    A: Snowflake maintains the view automatically on base-table change. MongoDB's \texttt{\$merge} pattern is maintained by your refresh job, so you own freshness, idempotency, and reconciliation.

    \item \textbf{Q: Why the lookback window in incremental refresh?}\\
    A: Frames reach backwards in time; a late event changes old frames. Refreshing watermark-minus-window-length re-computes every frame the late data could touch.

    \item \textbf{Q: What does \texttt{\$derivative} actually compute?}\\
    A: The slope between the first and last points of the frame, per unit time---a rate of change, not a point-to-point diff.
\end{enumerate}

\section{Further Reading}
The \texttt{\$setWindowFields} reference defines partition, frame, and operator semantics \cite{mongodbSetWindowFieldsDocs}; \texttt{\$merge} output behavior is in the aggregation stage documentation \cite{mongodbMergeDocs,mongodbAggregationDocs}. Kleppmann's chapters on batch and stream processing frame the materialization trade-offs \cite{kleppmann2017ddia}, and the Apache Spark project provides the batch-window comparison point used in the stretch goals \cite{apacheSparkRepo}.

\section*{Key Takeaways}
\begin{itemize}
    \item Windows compute per-document aggregates over ordered frames; \texttt{\$group} collapses; choose per question, not per habit.
    \item Deterministic sort keys make windows reproducible; ties without tiebreakers make them flaky.
    \item Accumulator semantics depend on context: \texttt{\$group}, \texttt{\$project}, and \texttt{\$setWindowFields} are three different contracts.
    \item \texttt{\$setWindowFields} is a blocking stage: filter first, exploit matching index order, respect the 100 MB limit.
    \item \texttt{\$merge} + watermark + deterministic keys is the materialized-view idiom; add lookback for late data.
    \item Every materialized view needs a reconcile job; freshness you cannot prove is freshness you do not have.
\end{itemize}

\section{Exercises}
\begin{enumerate}
    \item \textbf{(Conceptual)} Write the frame specification for ``previous 5 transactions plus current'' and for ``previous calendar week.''
    \item \textbf{(Conceptual)} Explain why a global leaderboard rank cannot be partitioned by region, and sketch the two-level alternative.
    \item \textbf{(Conceptual)} When is \texttt{\$expMovingAvg} preferable to a simple moving average over a range frame?
    \item \textbf{(Hands-on)} Add a \texttt{\$derivative} of queue depth per minute to the lab dataset and chart the result.
    \item \textbf{(Hands-on)} Break the materialized view by merging with a non-deterministic key, observe duplicates, then fix and prove idempotency.
    \item \textbf{(Design)} Design the refresh architecture for a tenant-isolated SaaS metrics product: watermarks per tenant, late data, backfill after an outage.
    \item \textbf{(Design)} Decide, with justification, whether a 30-day moving-average dashboard should be a live pipeline, a materialized view, or a cached API response at 1k, 100k, and 10M daily requests.
    \item \textbf{(Interview)} ``The materialized view drifted from the base data over a month.'' Enumerate the likely causes and your diagnostic order.
\end{enumerate}

\section{Capstone Project: Windowed Analytics with Incremental Materialization}\label{sec:ch11-capstone}

\begin{capstonebox}
\textbf{Mission:} build the windowed analytics platform for the orders domain: a 100-million-document workload served by ESR-compliant compound indexes (Chapter 9), faceted one-scan dashboards (Chapter 10), and windowed rolling metrics materialized with \texttt{\$merge} (this chapter)---with every plan proven covered and every refresh provably idempotent.

\textbf{Estimated time:} 6--8 hours, dominated by data generation and index builds.

\textbf{What you will practice:}
\begin{itemize}
    \item Integrating index architecture, pipeline design, and windowed analytics into one system.
    \item Proving zero \texttt{FETCH}/\texttt{COLLSCAN} and zero disk spills under realistic scale.
    \item Operating incremental materialization with lookback, reconciliation, and data-quality tests.
\end{itemize}
\end{capstonebox}

\subsection*{Scenario}
The payments platform from earlier capstones must now serve a finance analytics product: rolling 30-day customer spend metrics, multi-faceted revenue dashboards, and an organizational rollup of teams to regions via \texttt{\$graphLookup}. The workload is one hundred million ledger documents. Constraints: all dashboard queries fully covered, no disk spilling in steady state, and the analytics workload must not touch the OLTP primary---run it on a dedicated analytics replica set member or a hidden analytics node.

\subsection*{Requirements}
\begin{itemize}
    \item Generate 100M ledger documents (or a documented, proportionally scaled subset) with account, region, timestamp, and amount fields.
    \item Design one ESR compound index per dominant pipeline entry shape; validate coverage with \texttt{explain("executionStats")} on every pipeline.
    \item Implement the full pipeline chain: early \texttt{\$match} $\to$ \texttt{\$lookup} $\to$ \texttt{\$facet} $\to$ \texttt{\$setWindowFields} $\to$ \texttt{\$merge}.
    \item Implement incremental refresh with watermark plus 31-day lookback, deterministic merge keys, and a nightly reconcile job comparing one full partition against a reference computation.
    \item Provide at least three automated data-quality tests (aggregation parity vs. reference SQL, cardinality checks, empty-input edge case) with a failing-case demonstration.
    \item Provide a monthly cost estimate for the analytics replica set and storage footprint, with two optimization decisions documented.
\end{itemize}

\subsection*{Reference Architecture}
\begin{figure}[H]
    \centering
    \resizebox{0.95\textwidth}{!}{%
    \begin{tikzpicture}[
        st/.style={stage, minimum width=2.05cm, font=\scriptsize\bfseries},
        store/.style={cylinder, shape aspect=0.25, draw=primary, thick, fill=primary!8, shape border rotate=90, align=center, minimum width=1.5cm, minimum height=1.15cm, font=\scriptsize\bfseries},
        ops/.style={rectangle, rounded corners, draw=orange!85!black, fill=orange!10, minimum width=2.3cm, minimum height=0.9cm, align=center, font=\scriptsize\bfseries},
        arrow/.style={-{Stealth[length=2.5mm]}, draw=primary, thick}
    ]
        \node[store] (ledger) at (0,0) {ledger\\(100M)};
        \node[st] (match) at (2.6,0) {\$match\\ESR index};
        \node[st] (lookup) at (5.1,0) {\$lookup};
        \node[st] (facet) at (7.4,0) {\$facet};
        \node[st] (win) at (9.8,0) {\$setWindow-\\Fields};
        \node[st] (merge) at (12.3,0) {\$merge};
        \node[store] (mv) at (14.6,0) {metrics\\views};
        \node[ops] (wm) at (5.1,-1.7) {watermark +\\lookback 31d};
        \node[ops] (rec) at (12.3,-1.7) {nightly\\reconcile};
        \draw[arrow] (ledger) -- (match);
        \draw[arrow] (match) -- (lookup);
        \draw[arrow] (lookup) -- (facet);
        \draw[arrow] (facet) -- (win);
        \draw[arrow] (win) -- (merge);
        \draw[arrow] (merge) -- (mv);
        \draw[arrow, dashed] (wm) -- (match);
        \draw[arrow, dashed] (rec) -- (mv);
    \end{tikzpicture}%
    }
    \caption{Windowed analytics engine: one indexed scan feeds faceted windowed computation into reconciled materialized views.}
    \label{fig:ch8-capstone}
\end{figure}

\subsection*{Implementation Steps}
\begin{enumerate}
    \item \textbf{Baseline.} Generate data; run the naive pipeline; record the \texttt{COLLSCAN} baseline for the report.
    \item \textbf{Indexes.} Build ESR indexes per entry shape; prove \texttt{IXSCAN}-only plans and zero \texttt{FETCH} on covered projections.
    \item \textbf{Pipeline.} Assemble the full stage chain; verify determinism with the sort tiebreaker; confirm no spills with \texttt{allowDiskUse} disabled in steady state.
    \item \textbf{Materialize.} Add watermark, lookback, deterministic keys, and \texttt{\$merge}; demonstrate a late-arrival correction.
    \item \textbf{Verify.} Run the three data-quality tests, the reconcile job, and the failing-case demonstrations; write the cost estimate.
\end{enumerate}

\subsection*{Deliverables}
\begin{itemize}
    \item Generator, index migration, pipeline scripts, refresh job, and reconcile job.
    \item Explain evidence pack: IXSCAN-only, zero spills, covered projections.
    \item Data-quality test suite with failing-case output, plus the cost-estimate document.
\end{itemize}

\subsection*{Acceptance Criteria}
\begin{itemize}
    \item Every production pipeline starts with an index-backed \texttt{\$match}; no \texttt{COLLSCAN} appears in any explain output.
    \item Windowed results match the reference implementation on sampled partitions.
    \item Incremental refresh is idempotent and absorbs late arrivals within the lookback window.
    \item The reconcile job detects an injected drift in a test run.
\end{itemize}

\subsection*{Stretch Goals}
\begin{itemize}
    \item Replace the scheduled refresh with change-stream-triggered incremental updates (previewing Chapter 16).
    \item Benchmark the same rolling metrics in DuckDB or Spark and write a one-page build-vs-buy analysis.
\end{itemize}
